% SMPdesign/criteria.tex
% mainfile: ../perfbook.tex
% SPDX-License-Identifier: CC-BY-SA-3.0

\section{设计准则}
\label{sec:SMPdesign:Design Criteria}
%
\epigraph{一磅的学问需要十磅的常识来运用。}
	 {波斯谚语}

想要获取最佳的性能和可扩展性，简单的办法是不断尝试，
直到你的程序和最优实现水平相当。
可是如果你的代码不是短短数行，
可能的并行程序空间是如此巨大，
如何能在浩如烟海的写法中找到最优实现？
另外，"最佳的并行程序"到底是什么？
\cref{sec:intro:Parallel Programming Goals}
给出了三个并行编程目标：
\IX{performance}、\IX{productivity}和\IX{generality}，
而最佳性能常常要付出生产率和通用性的代价。
如果不在设计时就将这些选择考虑进去，
就很难在限定的时间内开发出性能良好的并行程序。

但除此以外，还需要更详细的设计准则来指导实际的设计，
这是本节要承担的任务。
在真实世界中，这些准则经常在某种程度上相互冲突，
这需要设计者小心权衡得失。

因此，这些准则可以被视为设计中的"阻力"，
而对这些阻力进行恰当权衡就被称为"设计模式"~\cite{Alexander79,GOF95}。

达到三个并行编程目标的设计准则是加速比、
竞争、开销、读写比和复杂性：
\begin{description}
\item[加速比：]  如\cref{sec:intro:Parallel Programming Goals}所述，
	提高性能是花费如此多时间和精力
	进行并行化的主要原因。
	加速比的定义是运行程序顺序版本所需时间
	与运行并行版本所需时间之比。
\item[竞争：]  如果对于一个并行程序来说，
	增加更多的CPU并不能让程序忙起来，
	多余的CPU会因竞争而无法做有用的工作。
	这可能是\IX{lock contention}、内存竞争，或者
	其他什么性能杀手。
\item[工作与同步比：]  一个单处理器、
	单\-/线程、不可抢占、且不可\-/中断\footnote{
		要么通过屏蔽中断，要么通过忽略中断。}
	版本的并行程序完全不需要任何同步原语。
	因此，这些原语消耗的任何时间
	（包括通信cache miss以及
	\IXh{message}{latency}、锁原语、原子指令
	和\IXpl{memory barrier}）
	都是对程序预期完成的有用工作没有直接贡献的开销。
	注意，重要的度量是同步开销
	与\IX{critical section}中代码开销之间的关系，
	较大的临界区能够容忍更大的同步开销。
	工作与同步比与同步效率的概念相关。
\item[读写比：]  对于极少更新的数据结构，
	通常可以采用"复制"而非"分区"，
	而且可以用非对称的同步原语来保护，
	以提高写者同步开销为代价来降低读者的同步开销，
	从而降低总体同步开销。
	对于频繁更新的数据结构也可以进行相应的优化，
	如\cref{chp:Counting}中所讨论的。
\item[复杂性：]  并行程序比等价的顺序程序更复杂，
	因为并行程序要比顺序程序维护更多的状态，
	尽管具有规则结构的大状态空间在某些情况下
	可以容易地被理解。
	并行程序员必须在这个更大的状态空间中
	考虑同步原语、消息传递、锁设计、
	临界区识别和死锁（deadlock）等诸多问题。

	更大的复杂性通常转化为
	更高的开发和维护成本。
	因此，预算约束可以限制对现有程序所做的
	修改的范围和类型，因为对原有程序的性能
	加速需要消耗相当的时间和精力。
	更糟糕的是，增加的复杂性甚至可能\emph{降低}
	性能和可扩展性。

	进一步说，超过某一点后，
	可能存在比并行化更廉价、更有效的
	潜在顺序优化。
	如\cref{sec:intro:Performance}所述，
	并行化只是众多性能优化之一，
	而且它最适用于CPU瓶颈。
\end{description}
这些准则将共同作用以强制限定最大加速比。
前三个准则是深度相关的，因此
本节的其余部分分析这些
相互关系。\footnote{
	一个真实世界的并行系统将受到许多额外的设计准则的约束，
	如数据结构布局、内存大小、内存层次延迟、
	带宽限制和I/O问题。}

注意，这些准则也可能作为需求规格说明的一部分出现，
此外它们是从\cpageref{sec:SMPdesign:Problems Quality Assessment}
总结并发算法质量这一问题的一个解决方案。
例如，加速比可以作为相对期望目标
（"越快越好"）
或作为工作负载的绝对要求（"系统
必须支持每秒至少1,000,000次web访问"）。
经典设计模式语言将相对期望目标描述为力，
将绝对要求描述为上下文。

理解这些设计准则之间的关系对于
识别并行程序的适当设计权衡取舍非常有帮助。
\begin{enumerate}
\item	程序在独占锁临界区上所花的时间越少，
	潜在的加速比就越大。
	这是\IXr{Amdahl's Law}~\cite{GeneAmdahl1967AmdahlsLaw}的结果，
	因为在给定时刻只有一个CPU可以在给定的
	独占锁临界区内执行。

	更确切地说，程序花在某个独占临界区中的时间比例
	必须大大小于CPU数的倒数，因为这样增加CPU数目才能达到事实上的加速。
	例如，除非程序在最受限的独占锁临界区上花费的时间
	远少于十分之一，否则程序不会扩展到10个~CPU。
\item	因竞争所浪费的大量CPU时间或挂钟时间，
	本来可用于提高加速比。实际的加速比应该尽可能接近可用CPU数目。
	CPU数量与实际加速比之间的差距越大，
	CPU的使用效率就越低。
	同样，期望的效率越高，
	可以继续提升的加速比就越小。
\item	如果使用的同步原语相较于它们保护的临界区
	开销太大，那么提高加速比的最佳方法
	是减少调用这些原语的次数。
	这可以通过批处理进入临界区、
	使用数据所有权（data ownership）（见\cref{chp:Data Ownership}）、
	使用非对称原语
	（见\cref{chp:Deferred Processing}）、
	或使用粗粒度设计如\IXh{code}{locking}来实现。
\item	如果临界区相较于保护它们的原语开销太大，
	那么提高加速比的最佳方法是通过转向
	读写锁、\IXh{data}{locking}、非对称原语
	或数据所有权来增加并行性。
\item	如果临界区相较于保护它们的原语开销太大，
	并且被保护的数据结构读多于写，
	那么增加并行性的最佳方法是转向
	读写锁或非对称原语。
\item	各种提高SMP性能的改动，例如
	减少锁竞争（lock contention），也会改善实时
	延迟~\cite{PaulMcKenney2005h}。
\end{enumerate}

\QuickQuiz{
	临界区的所有这些问题是否意味着
	我们应该始终使用
	非阻塞同步~\cite{MauriceHerlihy90a}，
	因为它没有临界区？
}\QuickQuizAnswer{
	虽然非阻塞同步在某些情况下非常有用，
	但它不是万能药，如
	\cref{sec:advsync:Non-Blocking Synchronization}中所讨论的。
	而且，正如Josh Triplett所指出的，
	非阻塞同步确实有临界区。
	例如，在基于compare-and-swap操作的
	非阻塞算法中，从初始加载开始到
	compare-and-swap的代码
	类似于基于锁的临界区。
}\QuickQuizEnd

值得重申的是，竞争有许多伪装，包括
锁竞争、内存竞争、cache溢出、\IX{thermal throttling}
等等。
本章主要关注锁竞争和内存竞争。
